% ABOUTME: Feature engineering section of the IEEE report (issue #4).
% ABOUTME: Rolling statistics, least squares slope, regime normalisation, health index, ablation.
\section{Degradation Features}
\label{sec:features}

Each FD001 row is one cycle of one engine: the cycle number, three
operating settings and 21 sensor readings. A single reading is noisy and
says little about how far an engine has degraded, so we add features that
summarise each sensor's recent history. All features are computed by
\texttt{make\_features} in \texttt{src/rul/features.py}, which returns one
row per input row with the original index. Every window is computed within
a single engine (grouped by \texttt{unit} and ordered by cycle), so no
feature ever mixes readings from two engines.

\subsection{Cycle}
The cycle number $t$ is kept unchanged as a feature. It is the time since
the engine entered service in the dataset, so on its own it is a weak prior
on wear. It is not scaled, since tree models are invariant to monotone
rescaling and downstream code uses the raw cycle.

\subsection{Operating regime normalisation}
Sensor levels depend on the flight condition (altitude, Mach number and
throttle angle) as much as on wear. We standardise the three settings and
cluster them with $k$-means (seed 42) into $K$ operating regimes. For each
regime $r$ and sensor $j$ we store the mean $\mu_{rj}$ and standard
deviation $\sigma_{rj}$ over the training rows in that regime, and
normalise
\begin{equation}
z_{t,j} = \frac{x_{t,j} - \mu_{r(t),j}}{\sigma_{r(t),j}},
\end{equation}
where $r(t)$ is the regime nearest to the row's settings. FD001 runs in one
regime ($K = 1$), so this is an ordinary z-score; with $K = 6$ the same code
removes the regime offsets of FD002 and FD004, leaving variation due to
wear. Sensors that are constant in FD001 (s\_1, s\_5, s\_10, s\_16, s\_18,
s\_19) have $\sigma = 0$ up to floating point noise; they are only centred,
so their features are zero rather than amplified noise.

The cluster centres and $(\mu, \sigma)$ are estimated on training engines
only and stored in a scaler object that is reused, unchanged, on
validation and test engines. If they were re-estimated on the validation
engines, the validation features would be scaled by statistics that
include those engines' own failure trajectories, and the cross validated
error would be optimistic. In cross validation the scaler is refitted in
every fold on that fold's training engines.

\subsection{Rolling mean and standard deviation}
For window length $w \in \{5, 10, 20\}$ and the set $W_t$ of the last
$\min(w, t)$ cycles of the same engine, with $n = |W_t|$,
\begin{equation}
m^{(w)}_{t,j} = \frac{1}{n}\sum_{s \in W_t} z_{s,j}, \qquad
s^{(w)}_{t,j} = \sqrt{\frac{1}{n - 1}\sum_{s \in W_t}
\bigl(z_{s,j} - m^{(w)}_{t,j}\bigr)^2}.
\end{equation}
The mean filters measurement noise and estimates the current level of
degradation; the standard deviation measures how erratic the sensor has
become, which tends to grow as components wear.

\paragraph{Early cycles} The first $w - 1$ cycles of an engine have fewer
than $w$ past readings. Rather than drop these rows or leave missing values,
the window is truncated to the cycles seen so far (\texttt{min\_periods=1}).
At cycle 1 the mean equals the reading and the standard deviation, undefined
for one value, is set to 0. These features are therefore noisier for the
first $w - 1$ cycles of each engine. Those cycles are at least 109 cycles
from failure in FD001, where the target is at or near the cap of 125, so
this costs little accuracy.

\subsection{Rolling least squares slope}
The slope is the ordinary least squares gradient of $z_{s,j}$ against cycle
$s$ over the same window. With $n$ points and sums taken over $W_t$,
\begin{equation}
b^{(w)}_{t,j} =
\frac{n \sum s\, z_{s,j} - \sum s \sum z_{s,j}}
     {n \sum s^2 - \bigl(\sum s\bigr)^2}.
\end{equation}
All five sums are rolling sums, so every slope is computed in closed form
in one pass instead of fitting a regression per row. For $n = 1$ the
denominator is zero and we define $b = 0$: a single point has no trend.

The slope is the local rate of degradation, in standard deviations per
cycle. The rolling mean says how worn an engine is; the slope says how fast
it is wearing. In FD001 the sensors drift slowly for most of an engine's
life and then accelerate as the fault develops, so a steep slope marks an
engine in the final phase of its life even when its sensor levels overlap
with those of a healthy engine of a different build. Because the slope uses
differences within one engine, it is also insensitive to the engine to
engine offsets in initial wear that make raw levels hard to compare. On a
window of 2 or 3 points the slope is dominated by noise, which is the main
cost of the truncated early windows.

\subsection{Health index}
The health index (HI) combines all normalised sensors into one number
between 0 and 1. We fit a linear regression on training engines,
\begin{equation}
h_t = 1 - \frac{t - 1}{T - 1}, \qquad
\widehat{\mathrm{HI}}_t = \mathrm{clip}\bigl(\beta_0 + \boldsymbol\beta^\top
\mathbf{z}_t,\ 0,\ 1\bigr),
\end{equation}
where $T$ is the engine's failure cycle, so the target $h$ falls linearly
from 1 at the first cycle to 0 at failure. At prediction time no $T$ is
needed: the fitted coefficients map any engine's sensors to an HI. We chose
regression over the first principal component because PCA finds the
direction of greatest variance, which need not be the direction of wear,
and its sign and scale are arbitrary; the regression picks the sensor
combination that tracks degradation directly and has a fixed 0 to 1 scale.
Its limitation is the linear target: real wear is flat for most of a life
and then accelerates, so the HI starts near 0.8 rather than 1 and only
falls clearly in the last third of life. On FD001 the HI has a correlation
of 0.86 with the capped RUL, higher than any single sensor (at most 0.80 in
absolute value for a 20 cycle rolling mean).

\subsection{Ablation}
Table~\ref{tab:feature-ablation} adds the feature families one at a time
to a Random Forest (default hyperparameters, seed 42), scored with five
fold cross validation grouped by engine on every validation cycle.

\begin{table}[t]
\centering
\caption{Feature ablation, 5 fold engine cross validation on FD001 training
engines (mean over folds, RUL capped at 125).}
\label{tab:feature-ablation}
\begin{tabular}{lrrrr}
\hline
Features & $n$ & RMSE & MAE & Late RMSE \\
\hline
Cycle and raw sensors      &  22 & 16.81 & 11.53 & 15.56 \\
+ rolling mean and std     & 148 & 16.35 & 11.02 & 13.87 \\
+ rolling slope            & 211 & 14.78 & 10.02 & 11.35 \\
+ health index             & 212 & 14.72 &  9.94 & 11.22 \\
\hline
\end{tabular}
\end{table}

Smoothing helps a little: a forest already averages over many noisy rows.
The slope gives the largest gain, about 1.6 RMSE and 2.5 late stage RMSE,
well beyond the fold to fold standard deviation of 0.5 to 0.8. The HI adds
a further small gain within that noise, since it is a linear combination of
sensors the forest already sees; its main value is as one interpretable
degradation signal.

\subsection{Limitation of fixed windows}
Every feature here looks back at most 20 cycles, while FD001 engines run
for 128 to 362 cycles. A fixed window cannot tell whether a sensor has been
drifting for 150 cycles or has just started to move, cannot see when the
drift began to accelerate, and must use one window length for engines that
degrade at very different rates. Lengthening the window trades this
against slower response and longer truncated starts, and the choice of
windows (5, 10, 20) is itself a hand-tuned hyperparameter. These are the
gaps that motivate Phase 2: sequence models such as LSTMs or temporal
convolutions take the whole trajectory as input and learn which part of
the history matters, rather than relying on fixed hand-made summaries.
